\section{Ordered fields and separation on the two arcs}
\label{ord:section}

The packing bound limits the number of secondary proposals that can
affect any one small ball. We now supply the scalar fields used by those
proposals and by
their final separation test. We use the circular order and the
noncrossing links of the column model.
All fields in this section are defined on a whole band $X_I$, including
its separate components. A bound for vertical Lipschitz constants,
together with compatibility at switches, implies the same bound for
the intrinsic distance $d_I$ wherever that distance is finite.

This section provides two tools.  First, a point-distance function has a
supporting field that is monotone toward a specified target along both
circular arcs and retains the exact distance at that target.  Second,
short paths on which the field rises and falls in a sufficiently thin band
force a distant pair with nearly equal field values onto opposite arcs.  A
signed transverse field will then distinguish the pair.

\subsection{Ordered supports of point distances}

\begin{definition}\label{ord:definition}
Let $a,b$ be two distinct column indices. A real field $s$ on $X_I$ is
\emph{ordered from $a$ to $b$}, with slope bound $M$, if it respects
every switch, is $M$-Lipschitz on every vertical column, and, at each
height $H\in I$, is nondecreasing on each of the two circular arcs
oriented from $a$ to $b$. In particular, the value on column $a$ is a
global minimum at that height, and the value on column $b$ is a global
maximum. We call the two open arcs the positive and negative arcs,
with either assignment of these names.
\end{definition}

Adding a constant to an ordered field preserves its order and slope
bound. Multiplication by a positive constant multiplies the slope
bound by that constant. Its negative is ordered after interchanging
the two boundary columns. A field depending only on height is
ordered for any choice of distinct boundary columns.

The supporting field below is obtained by taking running minima toward
the target column on each arc.  Connected distance sublevels are what make
this operation compatible with the switches.

\begin{lemma}[Ordered support of a point distance]\label{ord:support}
Suppose that $k,w$ belong to one component of $X_I$. There is an
ordered field $L$ of slope bound $1$ such that
\[
 L(k)=0,\qquad L(w)=d_I(k,w),\qquad
 L(v)\le d_I(k,v)\quad(v\in X_I),
\]
where the last inequality uses the convention that distance to a
different component is infinite.
\end{lemma}

\begin{proof}
Put $D=d_I(k,w)$. If $D=0$, the constant zero field suffices. If $k,w$
have representatives on the same column, then
$D=|h(w)-h(k)|$ and the field
\[
 L(v)=\min\{|h(v)-h(k)|,D\}
\]
suffices, with any two distinct indices as boundary columns: vertical travel attains this distance on that column, and
every path has length at least its absolute height change.

Otherwise choose a representative column $a$ of $k$ and a different
representative column $b$ of $w$. Write
\[
 f_i(H)=\min\{d_I(k,(i,H)),D\},\qquad H\in I,
\]
where $f_i=D$ on any component not containing $k$. Every $f_i$ is
$1$-Lipschitz, and the values respect switches. Moreover,
\[
 f_a(H)=\min\{|H-h(k)|,D\}=\min_i f_i(H).
\tag{\ref{ord:support}.1}\label{ord:root-minimum}
\]

For fixed $H$ and $t\le D$, every nonempty strict sublevel
\[
 S_t(H)=\{i:f_i(H)<t\}
\]
is connected in the active column graph and contains $a$. To see
this, let $i\in S_t(H)$ and take a path from $k$ to $(i,H)$ of length
less than $t$. If the path visits column $z$ at height $T$, its
remaining length is at least $|T-H|$. Its prefix followed by vertical
travel on $z$ to height $H$ therefore has length less than $t$.
Thus every visited column belongs to $S_t(H)$, and projection of
these paths to column indices connects $S_t(H)$ to $a$.

Consider a link with distinct endpoints $p,q$ at height $H$, and put
$c=f_p(H)=f_q(H)$. If neither endpoint is $a$, let $[p,q]_a$ be the
closed circular interval from $p$ to $q$ that does not contain $a$.
Then
\[
 f_z(H)\ge c\qquad(z\in[p,q]_a).
\tag{\ref{ord:support}.2}\label{ord:collision}
\]
Indeed, a strict-sublevel path from any column $z$ with $f_z(H)<c$
to $a$ avoids $p,q$, whereas Lemma~\ref{geo:chord-separation}
forbids such a path when $z$ is strictly inside $[p,q]_a$.
If one endpoint is $a$, both endpoint values equal the global
minimum in \eqref{ord:root-minimum}.

Cut the circular order at $a$. For $i\ne a$, let $[i,b]_a$ denote
the closed interval between $i$ and $b$ in the resulting linear
order, and define
\[
 L(i,H)=\min_{z\in[i,b]_a} f_z(H),\qquad
 L(a,H)=f_a(H).
\tag{\ref{ord:support}.3}\label{ord:projection}
\]
These are finite minima of $1$-Lipschitz functions. Along either
arc from $a$ to $b$ the intervals in \eqref{ord:projection} shrink,
so their minima are nondecreasing; the initial root value is a
lower bound for all of them by \eqref{ord:root-minimum}.

We check switches explicitly. If a link meets $a$, its other
endpoint has the global minimum value, and so does its value of
$L$. Otherwise order its endpoints as $p<q$ in the order cut at
$a$. By \eqref{ord:collision} the minimum of $f$ on $[p,q]_a$ is
their common value $c$. If $b$ lies between $p$ and $q$, both
minima defining $L(p,H),L(q,H)$ equal $c$. If $b$ lies outside
that interval, the longer of the two minimizing intervals is the
union of the shorter one and $[p,q]_a$. The minimum on the
shorter interval is at most $c$, because it contains a link
endpoint. The two minima are again equal. Thus $L$ respects
every switch and is an ordered field on the quotient.

Finally, $L(i,H)\le f_i(H)\le d_I(k,(i,H))$. At $k$, its
root value is zero. At $w$, the minimizing interval is the
singleton $\{b\}$ and its value is $D$.
\end{proof}

\subsection{A transverse coordinate}

At a fixed height the ordered field increases from $a$ to $b$ along both
arcs. Equal values can therefore occur on opposite arcs. The next field
records the side and vanishes wherever a switch can change sides.

\begin{definition}\label{ord:transverse-definition}
Let $s$ be ordered from $a$ to $b$ with slope bound $4$. In
$\R^2$ define the set
\[
 \begin{split}
 Z_s={}&\{(s(a,H),H),(s(b,H),H):H\in I\}\\
 &{}\cup\{(s(p,H),H): H\in I\cap S_{pq},\
     p\text{ and }q\text{ lie on opposite open arcs}\}.
 \end{split}
\]
For a representative $(i,H)$ define
\[
 t_s(i,H)=
 \begin{cases}
  \dist_\infty((s(i,H),H),Z_s),&i\text{ on the positive arc},\\
  -\dist_\infty((s(i,H),H),Z_s),&i\text{ on the negative arc},\\
  0,&i\in\{a,b\}.
 \end{cases}
\]
Here $\dist_\infty$ is distance for the maximum norm on $\R^2$.
\end{definition}

\begin{lemma}\label{ord:transverse-lipschitz}
The field $t_s$ is well defined on $X_I$, respects all switches,
and is $4$-Lipschitz for $d_I$.
\end{lemma}

\begin{proof}
The image of $H\mapsto(s(i,H),H)$ has Lipschitz constant at most
$4$ for the maximum norm, and distance to a fixed nonempty set
is $1$-Lipschitz. Thus $t_s$ is vertically $4$-Lipschitz on
each open-arc column; it is identically zero on each boundary
column. At a switch between columns on the same open arc, its
values agree because both $s$ and height agree. At a switch
between opposite open arcs the common coordinate image belongs
to $Z_s$, so both values vanish. At a switch involving a boundary
column, its image belongs to $Z_s$ for the same reason. These
checks also prove independence of representatives. Summing along
any finite path of vertical motions and switches proves the
$d_I$ bound.
\end{proof}

\subsection{Field excursions force opposite arcs}

The next observation records the crossing argument with its height
error. A path here can be a subpath of a concatenation; saying
that it stays within length $A$ of a point means that every point
on it can be joined to that point by a path of length at most $A$
in $X_I$. This condition holds, for example, for any
subpath of the union of two paths of length at most $A$ issuing
from that point.

\begin{lemma}[Overlapping traversals on one arc]\label{ord:overlap}
Let $J\subset I$ have width $\eta$. Suppose that $s$ is ordered
with slope bound $4$, and that paths $\Gamma_x,\Gamma_y$ in
$X_J$ use only columns of the same open arc. Assume that
$\Gamma_x$ stays within length $A_x$ of $x$ and $\Gamma_y$
stays within length $A_y$ of $y$. If the range of $s$ on each
path contains an interval $[u,v]$ with
\[
 v-u>32\eta,
\]
then
\[
 d(x,y)\le A_x+A_y+2\eta.
\]
\end{lemma}

\begin{proof}
Restrict and orient each path so that its endpoint values are
$u$ and $v$. Fix $H_*\in J$ and write $q_i=s(i,H_*)$.
The values $q_i$ are nondecreasing in the linear order on this
open arc, and
\[
 |s(i,H)-q_i|\le4\eta\qquad(H\in J).
\tag{\ref{ord:overlap}.1}\label{ord:profile-error}
\]
Choose a point on $\Gamma_x$ of value $c=(u+v)/2$, with
representative column $i$. Then $|q_i-c|\le4\eta$.
The initial column of $\Gamma_y$ has reference value at most
$u+4\eta<q_i$, and its terminal column has reference value
at least $v-4\eta>q_i$. Its column word therefore starts
strictly before $i$ and ends strictly after $i$.

If $\Gamma_y$ visits column $i$, connect the two visits on
that column vertically at cost at most $\eta$, giving the
claim. Otherwise Lemma~\ref{geo:skipping} supplies a switch
on $\Gamma_y$ with endpoint columns $p<i<q$. If its height
is $H_0\in J$, then
\[
 |q_p-q_q|\le8\eta,
 \qquad q_p\le q_i\le q_q.
\]
In particular both endpoint reference values differ from
$q_i$ by at most $8\eta$. The two endpoints of $\Gamma_x$
have reference values respectively at most $u+4\eta$ and
at least $v-4\eta$. Since $(v-u)/2>16\eta$, these lie
strictly below $q_p$ and strictly above $q_q$, respectively.
Thus the middle point on column $i$ is strictly inside the
interval bounded by the link $pq$, whereas the endpoints
of $\Gamma_x$ are strictly outside it.
Lemma~\ref{geo:chord-separation} implies that $\Gamma_x$
visits $p$ or $q$. Connect this visit vertically to height
$H_0$ and, if needed, use the zero-cost switch $pq$.
This joins a point of $\Gamma_x$ to a point of $\Gamma_y$
at cost at most $2\eta$ (in fact $\eta$ suffices when
using the visited endpoint of that switch). The assumed
connections to $x,y$ complete the proof.
\end{proof}

\begin{lemma}[Quantitative separation on the two arcs]\label{ord:two-side}
Set $R=10^{-8}$. Let $s$ be an ordered field of slope bound
$4$ on $X_I$. Suppose that $x,y\in X_I$ satisfy
\[
 1\le d(x,y)\le2,\qquad d_I(x,y)=d(x,y),
 \qquad |s(x)-s(y)|\le10^{-3}R.
\]
Let $J\subset I$ be a common height interval of width
\[
 \eta\le10^{-6}R
\]
containing $x,y$. For each $z\in\{x,y\}$, suppose that
there are two paths $\gamma_z^-,\gamma_z^+$ in $X_J$,
starting at $z$, of length at most $2R$, whose terminal
points $z^-,z^+$ satisfy
\[
 s(z^-)<s(z)-0.4R,\qquad s(z^+)>s(z)+0.4R.
\]
Then $t_s(x),t_s(y)$ have opposite signs and
\[
 0.01R\le |t_s(x)|,|t_s(y)|\le8.
\]
\end{lemma}

\begin{proof}
Write $s_0=s(x)$ and $H_x=h(x)$. We first prove that $Z_s$
does not meet the rectangle
\[
 \mathcal Q=
 [s_0-0.2R,s_0+0.2R]\times
 [H_x-0.02R,H_x+0.02R].
\tag{\ref{ord:two-side}.1}\label{ord:empty-rectangle}
\]
Suppose otherwise that $(c,H_0)\in Z_s\cap\mathcal Q$.
At the reference height $H_0$, the endpoint columns of
each low path have field value strictly below $c$, and
the endpoint columns of each high path have value strictly
above $c$. Indeed their terminal values are respectively
below $s_0-0.399R$ and above $s_0+0.399R$, and moving
their values to height $H_0$ changes them by at most
\[
 4(0.02R+\eta)\le0.080004R.
\]
These inequalities leave a strict margin from the interval
$[s_0-0.2R,s_0+0.2R]$.

If $(c,H_0)$ is the image of a boundary column, these
strict inequalities contradict that column's extremality.
Otherwise it is the image of a crosslink with one endpoint
on each open arc. At height $H_0$, monotonicity puts every
column of value less than $c$ on the root side of this
crosslink and every column of value greater than $c$ on
its opposite side. The projected concatenation
$(\gamma_z^-)^{-1}\gamma_z^+$ must therefore visit an
endpoint of the crosslink by
Lemma~\ref{geo:chord-separation}. Each such visit is
reachable from $z$ in length at most $2R$, since it lies
on one of the two original paths. Move vertically from
the two visits to height $H_0$ and use the crosslink if
needed. This produces an $x$--$y$ path of length at most
\[
 4R+2(0.02R+\eta)<1,
\]
contradicting $d(x,y)\ge1$. This proves
\eqref{ord:empty-rectangle}.

The maximum-norm ball of radius $0.01R$ around either
$(s(x),h(x))$ or $(s(y),h(y))$ is contained in
$\mathcal Q$. Hence
\[
 |t_s(x)|,|t_s(y)|\ge0.01R.
\tag{\ref{ord:two-side}.2}\label{ord:positive-magnitude}
\]
In particular both centers have representatives on open
arcs, with a well-defined choice of arc.

Consider the value strip
\[
 \mathcal S=\{v\in X_J: s_0-0.18R\le s(v)
                                      \le s_0+0.18R\}.
\]
No representative of a point of $\mathcal S$ is on a
boundary column, and no switch in $\mathcal S$ joins
opposite open arcs. Either event would give an image in
$Z_s\cap\mathcal Q$, because its height belongs to $J$
and its value is within $0.18R$ of $s_0$. Thus any path
contained in $\mathcal S$ remains on a single open arc.

For $z=x,y$, the concatenation
$(\gamma_z^-)^{-1}\gamma_z^+$ starts below the strip
and ends above it. The scalar field is continuous along
this path, including at switches. Take its first visit
to the upper strip boundary, and its last visit to the
lower strip boundary preceding that visit. The intervening
subpath $\Gamma_z$ stays in $\mathcal S$, traverses its
full value interval, and lies on one open arc. Every
point on it is reachable from $z$ in length at most $2R$.
If $\Gamma_x,\Gamma_y$ were on the same arc,
Lemma~\ref{ord:overlap}, with common interval of length
$0.36R>32\eta$, would give
\[
 d(x,y)\le4R+2\eta<1.
\]
Consequently their two arcs are opposite.

It remains to relate these full traversals to the centers.
From either center $z$, follow either one of its two
original paths until its first exit from $\mathcal S$.
The resulting initial subpath $\Pi_z$ lies on the
center's arc and its value range contains the interval
between $s(z)$ and the strip boundary it first reaches.
This interval has length at least $0.179R$. All its
points are reachable from $z$ in length at most $2R$.
If the two centers had the same arc, exactly one of the opposite
full traversals $\Gamma_x,\Gamma_y$ would lie on that
arc. If it were $\Gamma_x$, compare it with $\Pi_y$; if
it were $\Gamma_y$, compare it with $\Pi_x$. In either case
the paths come from different centers and their value ranges overlap in
an interval
of length at least $0.179R>32\eta$. Another application
of Lemma~\ref{ord:overlap} gives the same contradiction.
Thus the centers have opposite arcs and, by
\eqref{ord:positive-magnitude}, $t_s(x),t_s(y)$ have
opposite signs.

Finally Lemma~\ref{ord:transverse-lipschitz} and the
assumed equality of band and unrestricted distances give
\[
 |t_s(x)|+|t_s(y)|
   =|t_s(x)-t_s(y)|
   \le4d_I(x,y)=4d(x,y)\le8.
\]
This proves both upper bounds.
\end{proof}


\subsection{Ordered charts in system \texorpdfstring{$B$}{B}}
\label{ord:charts}

Besides ordinary charts, system $B$ uses \emph{ordered charts}. Such a
chart has field
\[
 P=s+\lambda u,
\]
where $s$ is ordered with vertical Lipschitz constant at most $4$ and
$u$ has Lipschitz constant at most $K$ for $d_I$. At the creation of the
chart, $u=0$. The update in Section~\ref{sec:increment} changes $u$ while
preserving this bound. Our choice of $\lambda$ gives
$\Lip(P)\le4+\lambda K<5$ for finite $d_I$, as required by the chart
invariant of Section~\ref{sets:section}.

The band $I$, the ordered part $s$, its two boundary columns and choice
of positive arc, and the label functions $Z,\sigma$ remain fixed in original units for the
lifetime of an ordered identifier. In particular its transverse field
$t_s$ remains fixed. Its region can shrink as other charts replace parts
of it, and its auxiliary field $u$ changes through the prescribed
increments. Thus the current full field $P$ may change without changing
the identifier or the ordered data it records.
